% Part: sets-functions-relations
% Chapter: sets
% Section: subsets

\documentclass[../../../include/open-logic-section]{subfiles}

\begin{document}

\olfileid{sfr}{set}{sub}
\olsection{Deelverzamelingen en machtsverzamelingen}

\begin{explain}
We willen verzamelingen vaak vergelijken. Een duidelijke vraag is dan:
\emph{zit alles uit de ene verzameling ook in de andere?} Deze situatie is
belangrijk genoeg om er een nieuwe schrijfwijze voor in te voeren.
\end{explain}

\begin{defn}[Deelverzameling]
Als elk !!{element} van een verzameling $A$ ook !!a{element} van~$B$
is, heet $A$ een \emph{deelverzameling} van~$B$. We schrijven dan
$A \subseteq B$. Is $A$ geen deelverzameling van~$B$, dan schrijven
we $A \not\subseteq B$. Als $A \subseteq B$ maar $A \neq B$, schrijven
we $A \subsetneq B$. Dan heet $A$ een \emph{echte deelverzameling}
van $B$.
\end{defn}

\begin{ex}
Elke verzameling is een deelverzameling van zichzelf. Ook is
$\emptyset$ een deelverzameling van elke verzameling. De even
natuurlijke getallen vormen een deelverzameling van de natuurlijke
getallen. Verder geldt $\{ a, b \} \subseteq \{ a, b, c \}$. Maar
$\{ a, b, e \}$ is geen deelverzameling van $\{ a, b, c \}$.
\end{ex}

\begin{ex}
Het getal $2$ is een !!{element} van de verzameling van de gehele
getallen. De verzameling van de even getallen is een deelverzameling
van die verzameling. Toch kan een verzameling \emph{tegelijk} !!a{element}
en een deelverzameling van een andere verzameling zijn. Zo geldt
$\{0\} \in \{0, \{0\}\}$, maar ook
$\{0\} \subseteq \{0, \{0\}\}$.
\end{ex}

De regel van extensionaliteit zegt wanneer twee verzamelingen gelijk
zijn: $A = B$ precies als elk !!{element} van~$A$ ook !!a{element}
van~$B$ is, en omgekeerd. De eerste helft daarvan is precies wat
$A \subseteq B$ betekent: elk !!{element} van~$A$ is ook
!!a{element} van~$B$. Wisselen we $A$ en $B$ om, dan krijgen we
$B \subseteq A$ precies als elk !!{element} van~$B$ ook !!a{element}
van~$A$ is. Dat is de andere helft van ‘en omgekeerd’. Samen zeggen
die twee helften dat verzamelingen precies gelijk zijn als ze elk een
deelverzameling van de andere zijn.

\begin{prop}
$A = B$ precies als $A \subseteq B$ én $B \subseteq A$.
\end{prop}

Nu voeren we nog wat handige tekens in. Bij de uitleg van
$A$ als deelverzameling van~$B$ schreven we ‘elk !!{element} van~$A$
is \dots,’ en op de ‘$\dots$’ kwam ‘!!a{element} van $B$’. Zinnen met
deze \emph{vorm} komen zo vaak voor dat een korte schrijfwijze handig
is.

\begin{defn}\ollabel{forallxina}
$(\forall x \in A)\phi$ is een korte schrijfwijze voor $\forall x(x
\in A \lif \phi)$. Net zo is $(\exists x \in A)\phi$ een korte
schrijfwijze voor $\exists x(x \in A \land \phi)$.
\end{defn}

Met deze tekens zeggen we: $A \subseteq B$ precies als
$(\forall x \in A)x \in B$.

Nu kijken we naar de verzameling waarin alle deelverzamelingen van één
gegeven verzameling zitten.

\begin{defn}[Machtsverzameling]
Neem alle deelverzamelingen van een verzameling~$A$ bij elkaar. De
verzameling die je dan krijgt, heet de \emph{machtsverzameling van}~$A$.
We schrijven daarvoor $\Pow{A}$.
  \[
    \Pow{A} = \Setabs{B}{B \subseteq A} 
  \]
\end{defn}

\begin{ex}
Welke deelverzamelingen van $\{ a, b, c \}$ zijn er? Dat zijn:
$\emptyset$, $\{a \}$, $\{b\}$, $\{c\}$, $\{a, b\}$, $\{a, c\}$,
$\{b, c\}$ en $\{a, b, c\}$. Als we die allemaal in één verzameling
zetten, krijgen we $\Pow{\{a,b,c\}}$:
\[
\Pow{\{ a, b, c \}} = \{\emptyset, \{a \}, \{b\}, \{c\}, \{a, b\},
\{b, c\}, \{a, c\}, \{a, b, c\}\}
\]
\end{ex}

\begin{prob}
Schrijf alle deelverzamelingen van $\{a, b, c, d\}$ op.
\end{prob}

\begin{prob}
Laat zien dat als $A$ $n$ !!{element}s heeft, $\Pow{A}$ $2^n$
!!{element}s heeft.
\end{prob}
\end{document}
